% =============================================================================
% SECTION 1: INTRODUCTION
% =============================================================================
\section{Introduction}

% --- Two-component architecture ---
Modern malware that delivers arbitrary code typically consists of two components: a \textit{payload}, the shellcode that performs the intended action such as a reverse shell, and a \textit{loader}, the program that prepares, places, and executes the payload in memory.
The payload defines what the malware does; the loader defines how the payload reaches execution.

% --- Why loaders matter ---
Most evasion engineering concentrates in the loader.
Antivirus (AV) and Endpoint Detection and Response (EDR) products observe loader behavior more than payload content: they scan embedded data at rest, monitor memory allocation, and flag suspicious thread creation and code injection~\cite{mitre_execution}.
The payload itself is often generic and interchangeable.
This paper focuses on \textit{shellcode loaders}---programs that load and execute position-independent code in memory on Windows---rather than DLL loaders, script-based loaders, or fileless techniques targeting other platforms.
Consistent with this focus, we treat the shellcode payload as a controlled variable: all experiments use a single \texttt{msfvenom}-generated reverse-TCP shellcode as a fixed reference payload, so that observed detection differences can be attributed to loader choices rather than to payload content.

% --- Problem statement ---
Despite their importance, existing knowledge about loaders is fragmented.
Offensive techniques such as process injection~\cite{process_injection}, API unhooking~\cite{api_unhooking}, and direct syscall invocation~\cite{hellsgate} are documented in scattered blog posts, repositories, and conference talks, with no shared framework that relates them within a loader's execution flow.
On the defensive side, analysts often treat loaders as black boxes.
When a loader evades detection, it is hard to tell which stage was responsible---encrypted storage, direct syscalls, or callback-based execution---and each requires a different detection strategy.

% --- Brief gap motivation ---
Existing approaches address parts of this problem.
MITRE ATT\&CK~\cite{mitre_attack} catalogs adversary tactics and techniques but does not model the internal stages of a single loader binary.
Sandbox platforms such as Cuckoo~\cite{cuckoo} and ANY.RUN~\cite{anyrun} run samples as opaque inputs and do not support technique-level comparison across controlled variants.
Section~2 reviews related work in detail.

% --- Contributions ---
This paper makes three contributions:

\begin{enumerate}[leftmargin=*, itemsep=2pt]
    \item A six-stage pipeline model that decomposes a shellcode loader into sequential stages---Anti-Analysis, Storage, Allocation, Transformation, Writing, Execution---together with a notation $L_i.T_j$ for identifying techniques within each stage.

    \item An automated build-and-test platform that instantiates the model through compile-time technique selection, runs generated variants in isolated Windows virtual machines, and collects stage-resolved telemetry from Windows Defender and Sysmon.

    \item A stage-aware detection methodology that links each pipeline stage to candidate telemetry sources and rule formats, providing a structured framework for future detection-rule development; rule implementation and empirical validation are scoped as future work.
\end{enumerate}

The model is intended to be descriptive rather than prescriptive: its purpose is to express any shellcode loader---including those produced by existing operator frameworks---as a tuple $C$ of stage choices, so that loaders can be compared, classified, and studied in a common vocabulary.
Windows Defender serves throughout the paper as a \textit{detection reference} used to check that stage-level distinctions correspond to observable differences in loader behavior.
It is not a target of the study.

% --- Key findings preview ---
Preliminary experiments using Windows Defender as a detection reference indicate that stage-level changes produce distinct observable patterns: the Transformation stage dominates static-detection outcomes, while Execution techniques and API mode influence behavioral telemetry.
These observations support the decomposition's stage boundaries as capturing distinctions that matter to detection behavior---evidence for the model's classification basis, not a claim about evasion effectiveness.

% --- Paper organization ---
The rest of the paper is organized as follows.
Section~2 formalizes the shellcode/loader decomposition and reviews related work.
Section~3 presents the pipeline model and the platform.
Section~4 reports experimental results.
Section~5 discusses findings and limitations.
Section~6 concludes.
